\chapter{Conclusion}
\label{ch:conclusion}

The end of this chapter summarises the work completed by MedLift-3D, its limitations and future work, and introspects on the independent learning needed by the project.

\section{Summary}

MedLift-3D converts a single click on a low dose chest tomosynthesis image to a 3D measurement of the nodule, and provides the calibrated uncertainty measurement. We demonstrated via projection moments that DTS measures a nodule's mass in all views but its depth extent about 15 times worse, and built the system around this fact: a prism reconstruction around the click, a 17.4-million-parameter conditional diffusion prior that was consistent with the X-rays, 3D Gaussians that give mass in closed form, conformal intervals and calibrated size-category probabilities, a learned detector to suggest clicks, and a seen/unseen explanation of every result. A virtual clinical trial was conducted on LIDC-IDRI, and verified 18-step pipeline ran the trial with 583 patients. The components performed as desired: the prior raised the nodule-region PSNR from 11.88 dB to 15.23 dB on the held-out data, the detector achieved a CPM of 0.298, versus 0.017 for a training-free baseline, the conformal intervals had their desired coverage, the calibrated probabilities reduced the Brier score for $P(V\geq100)$ from 0.54 to 0.17, and the null-space analysis demonstrated that the X-rays measure nearly all of a nodule's mass and that the prior's unseen content correlates with the truth. The measurements did not meet the desired accuracy: Mass was not more accurate than Volume (H1); the whole method was not better than the physics-only reconstruction method (H2); MDT was not better than VDT with our method (H4); and the valid intervals were too large to be useful (H3). Calibrated probabilities (H5) were helpful at one of two thresholds, while the positive (unseen) prior skill from the model favors H6. The analysis revealed the failures to specific causes and fixes and demonstrated that it was not the case of better image quality is better measurement -- also a lesson for medical generative AI.

\section{Limitation}

\begin{itemize}
  \item \textbf{Simulation only.} There is no public chest DTS dataset with CT ground truth; the performance of the DTS on real scanners is unknown since our DTS is simulated using CT, monochromatic beam and no scatter.\item \textbf{Acquisition.} For example 15 views as opposed to the 60-74 views of clinical units (at equal dose) and a conservative \ang{30} sweep; the prior was trained and tested only at this geometry.\item \textbf{Detector sensitivity.} 35\% at 2 false positives per scan is not sufficient for unaided use - detector is an aid only.\item \textbf{Growth.} The growth of the virtual nodule is uniform whereas real nodules grow irregularly.\item \textbf{Ground truth.} LIDC labels vary between radiologists (consensus volume / mean reader volume 0.81), and only nodules with a diameter of 4--30 mm which were marked by at least 3 readers were used.\item \textbf{Population.} LIDC comprised primarily of US patients, and did not represent populations with tuberculosis scarring.\item \textbf{Measurement accuracy.} Volume and mass errors (MdAPE 75--100\%) are unsuitable for clinical practice and the whole method under-estimates considerably (bias $-91\%$) and fails in 22\% of nodules.\item \textbf{Uninformative intervals.} The intervals were valid and unbounded, and the ensemble spread was negatively correlated with the error due to failed measurements in the calibration set.\item \textbf{Training budget and size.} The prior continued to improve at the end of training, and its samples were more useful for training patients than for new ones (+7.00 vs +4.22 dB), and the detector also partially memorised its training patients.\item \textbf{Sample size.} The numbers of nodules used for measurements were 150 (76 calibration and 74 test) and 40 growth pairs and the difference between the methods is small and therefore not conclusive.\item \textbf{Patient thickness} was obtained from the CT body outline and for a real DTS unit it would have to be measured.
\end{itemize}

\section{Future Work}

\begin{enumerate}
  \item \textbf{Nodule-preserving prior:} condition the diffusion model on the click position, and oversample the pairs of nodules, so that faint nodules aren't replaced by normal lung.\item \textbf{Projection-domain mass:} use integration of the detector signal around the nodule's shadow, which is the null-space analysis' indicator of the signal being fully measured, as opposed to integration of a segmented reconstruction.\item \textbf{Failure-aware calibration:} Intervals are calibrated on measurable cases only, and additional outcome is ``not measurable'' with its own probability and a difficulty score is positively correlated with error.\item \textbf{Real DTS data:} collab with a hospital with a DTS unit to obtain paired DTS and CT scans, re-calibrate the conformal layer on these and measure the real performance.\item \textbf{Physics realism:} incorporate polychromatic spectra, scatter and the clinical protocols (60--74-view) in the simulator.\item \textbf{Longer and larger training:} train the prior longer with augmentation and test the larger models and train detector with more nodule-free negatives.\item \textbf{Real growth:} don't use virtual growth, instead use longitudinal CT pairs.\item \textbf{Reader study:} evaluate the accuracy of follow-up decision making, when using MedLift-3D in addition to DTS, compared to DTS-only use.\item \textbf{Deployment:} a Bangla interface, DICOM/PACS integration, an on-site GPU option and a quality system towards registration.\item \textbf{Other limited-angle problems:} breast tomosynthesis and dental cone-beam CT which have the commonality of the missing-wedge structure.
\end{enumerate}

\section{Independent Learning and Reflection}

This project involved much more knowledge and skill than was being taught in our course. In this section, we look back at that independent learning under five aspects: curiosity, self-initiative, independence, transfer of past learning, reflection on what we learned.

\subsection{Curiosity}
This was originally conceived as a generative-modelling concept but we were continually asking ourselves ``what can a \ang{30} X-ray sweep actually measure?'' Taking that question over one of the projection moments resulted in the mass-first design. Due to this curiosity, we wondered if the prior recovered or invented depth, and, in fact, the seen/unseen analysis came into existence and we were able to produce calibrated category probabilities only after studying conformal prediction.

\subsection{Self-initiative}
In addition to the mandatory scope we included a learned detector after the training-free detector, which mainly detected vessels, a fit check to demonstrate that no overfitting was occurring in the prior, conformal predictive systems with size-group calibration, explainability with a truth referenced test and a geometry sensitivity study to justify the 15-view choice, and an automatic HTML report that made all checks visible to examiners.

\subsection{Independence}
In addition to the classroom, we were able to learn about running a long unattended job on rented cloud machines, downloading and organizing 64 GB of medical images from TCIA, using pylidc to read radiologist annotations and serving a GPU application from a Kaggle notebook via an HTTPS tunnel when the traditional share link didn't work.

\subsection{Transfer of Past Learning}
A lot of the project is based on knowledge developed in previous courses which is applied to a problem none of the courses dealt with. Linear algebra was used to disentangle what the X-rays measured from what the prior supplies; calculus was used to give the moment argument that underpinned the design; probability and statistics became calibrated uncertainty; and software engineering was used to turn the design into a verified pipeline. \Cref{tab:transfer} shows the areas in which each application was made.

\begin{table}[!htbp]
\centering\small
\caption{Earlier coursework applied in MedLift-3D.}
\label{tab:transfer}
\begin{tabularx}{\linewidth}{>{\raggedright\arraybackslash}p{4.8cm} >{\raggedright\arraybackslash}X}
\toprule
Earlier course knowledge & Application in MedLift-3D \\
\midrule
Linear algebra (range, null space, least squares) & Seen/unseen split with CGLS; SIRT as an iterative least-squares solver \\
Calculus and moments & Projection-moment argument: mass is the 0th moment, depth extent the 2nd \\
Probability and statistics & Poisson noise, conformal prediction, Brier score, rank correlations, bootstrap confidence intervals \\
Numerical methods & Iterative solvers, numerical stability, mixed precision \\
Data structures and algorithms & Ray marching, candidate ranking, patient-level splitting \\
Artificial intelligence and machine learning & Diffusion training, U-Nets, overfitting checks, train/calibration/test discipline \\
Image processing and computer graphics & Projection geometry, 3D Gaussians, segmentation, saliency maps \\
Software engineering & Requirements and traceability, unit tests, version control, the verified pipeline \\
Computer networks & HTTPS tunnels and port forwarding for the application \\
\bottomrule
\end{tabularx}
\end{table}

\subsection{Reflection on Learning}
Our biggest lesson is that a large architecture with a decreasing training loss does not mean that one clinical measurement is correct. The restructuring of the work on measurable clinical numbers, explicit hypotheses and automatic checks had an impact on the way people think about engineering: if we can't verify a result, it isn't a result; if it is a negative result, it is more valuable if we tell the truth. We also learned to work within real constraints of time, money and data, to budget our computation as well as to embrace uncertainty as part of our answer instead of an embarrassment. It's not so much a model we take away from the project as these habits.
